\answer{OA17}
设总共有 $N$ 种，第 $j$ 种出现在 $r_j$ 个人的收藏中，其中 $1\le r_j\le4$。计数所有人持有的种类数，得到
\[
 \sum_{j=1}^N r_j=40.
\]
再把六对收藏的交集大小相加。第 $j$ 种被计入 $\binom{r_j}{2}$ 次，因此
\[
 \sum_{j=1}^N\binom{r_j}{2}\le36,
 \qquad \sum_{j=1}^N r_j^2\le112.
\]
Cauchy--Schwarz 不等式给出
\[
 40^2\le N\sum_{j=1}^N r_j^2\le112N,
 \qquad N\ge\frac{100}{7}.
\]
$N$ 是整数，故 $N\ge15$。

把明信片类型编号为 $1$ 至 $15$，取以下四个收藏：
\[
 \begin{aligned}
 A&=\{1,2,3,4,5,6,7,8,9,10\},\\
 B&=\{1,2,3,4,5,6,11,12,13,14\},\\
 C&=\{1,2,7,8,9,10,11,12,13,14\},\\
 D&=\{3,4,5,7,8,9,11,12,13,15\}.
 \end{aligned}
\]
每个集合都有十个元素，六个两两交集均恰有六个元素，并集包含全部十五种。因此最小值为 $\boxed{15}$。
